\subsection{非迟滞极限：静电反射与闭合损失谱}
\label{chap:feqo-eels}

完整 Fresnel 积分包含传播与倏逝两类场。为了检验归一化并分离材料极点和轨迹权重，下面在同一界面上取非迟滞极限。沿用界面 $z=0$、电子沿 $x$ 运动的坐标；它与电势计算中常用的法向 $x$、运动方向 $z$ 只是坐标互换。

\begin{proposition}[半空间的静电反射核]
设真空在 $z>0$，材料在 $z<0$，电子轨迹为 $(vt,0,b)$。以 $\ee^{-\ii\omega t}$ 为时间约定，并对界面方向作 Fourier 变换，记 $\kappa=(\kappa_x^2+\kappa_y^2)^{1/2}$。在贡献波数上，要求真空侧 $\omega/(c\kappa)\ll1$、材料侧 $|\sqrt\varepsilon|\omega/(c\kappa)\ll1$；两侧法向波数才可同时约化为静电衰减常数。忽略相应迟滞修正后，诱导电势为
\begin{equation}
\widetilde\varphi_{\rm ind}(z,\boldsymbol\kappa,\omega)
=\frac{q_{\mathrm e}}{2\varepsilon_0v\kappa}
r_{\rm es}(\omega)\ee^{-\kappa(z+b)}
2\pi\delta(\kappa_x-\omega/v),
\qquad r_{\rm es}=\frac{1-\varepsilon}{1+\varepsilon}.
\label{eq:feqo-planar-image-potential}
\end{equation}
这里 $r_{\rm es}$ 是电势反射系数；它与本章 TM 电场振幅在静电极限的符号相反。
\end{proposition}
\begin{proof}
电子电荷密度的时间变换为 $q_{\mathrm e}\delta(y)\delta(z-b)\ee^{\ii\omega x/v}/v$。界面 Fourier 变换后，Poisson 方程成为
\[
(\partial_z^2-\kappa^2)\widetilde\varphi=-\widetilde\rho/\varepsilon_0.
\]
点源导数跳跃给出自由项系数 $1/(2\kappa)$。源和界面之间的自由项正比于 $\ee^{\kappa z}$，反射项正比于 $r\ee^{-\kappa z}$，材料中的衰减项正比于 $t\ee^{\kappa z}$。

电势连续与无自由面电荷分别给
\[
1+r=t,\qquad 1-r=\varepsilon t.
\]
消去 $t$ 得 $r=r_{\rm es}$。当 $\varepsilon=1$ 时反射场消失；被动材料满足
\[
-\operatorname{Im}r_{\rm es}
=\frac{2\operatorname{Im}\varepsilon}{|1+\varepsilon|^2}\ge0.
\]
\end{proof}

\begin{theorem}[单位长度非迟滞损失谱]
在长直轨迹、线性响应和弱单次损失条件下，非迟滞概率密度为
\begin{align}
\frac1L\frac{\dd P}{\dd\omega}
&=\frac{q_{\mathrm e}^2}{2\pi\varepsilon_0\hbar v^2}
\int_{-\infty}^{\infty}\frac{\dd\kappa_y}{2\pi}
\frac{\ee^{-2b\sqrt{\kappa_y^2+(\omega/v)^2}}}
{\sqrt{\kappa_y^2+(\omega/v)^2}}
[-\operatorname{Im}r_{\rm es}],\label{eq:feqo-planar-eels}\\
&=\boxed{\frac{q_{\mathrm e}^2}{2\pi^2\varepsilon_0\hbar v^2}
K_0(2\omega b/v)[-\operatorname{Im}r_{\rm es}].}
\label{eq:feqo-planar-eels-k0}
\end{align}
其单位为 ${\rm s\,m^{-1}}$。近似要求实际贡献的波数满足 $\kappa\gg\max\{1,|\sqrt\varepsilon|\}\omega/c$，不能只用颗粒尺寸或“近场”一词替代该条件。在窄共振处还须检查被忽略修正相对反射分母 $1+\varepsilon$ 是否足够小；完整 Fresnel 积分负责处理该极限不均匀的区域。
\end{theorem}
\begin{proof}
将 $E_{{\rm ind},x}=-\ii\kappa_x\varphi_{\rm ind}$ 代入电子损失功，再使用 $\kappa_x=\omega/v$，得到第一式。它也可由完整 TM 核检验：$\alpha_1\to\kappa$、$r_p\to-r_{\rm es}$，TE 项相对消失，$\mu_0/k_0^2=1/(\varepsilon_0\omega^2)$，恰好恢复相同系数。

令 $\kappa_y=(\omega/v)\sinh u$，则
\[
\int_{-\infty}^{\infty}\frac{\dd\kappa_y}{2\pi}
\frac{\ee^{-2b\kappa}}\kappa
=\frac1\pi\int_0^\infty\ee^{-(2\omega b/v)\cosh u}\dd u
=\frac1\pi K_0(2\omega b/v).
\]
这给出第二式，并由被动性保证概率非负。
\end{proof}

\begin{workedexample}[Drude 峰位、线宽与距离权重]
取 $\varepsilon=1-\omega_p^2/[\omega(\omega+\ii\Gamma_D)]$，则
\[
-\operatorname{Im}r_{\rm es}
=\frac{2\Gamma_D\omega\omega_p^2}
{(2\omega^2-\omega_p^2)^2+4\Gamma_D^2\omega^2}.
\]
写 $\omega_s=\omega_p/\sqrt2$、$\delta=\omega-\omega_s$，在 $|\delta|,\Gamma_D\ll\omega_s$ 下，
\[
-\operatorname{Im}r_{\rm es}
\simeq\frac{\Gamma_D\omega_s}{4\delta^2+\Gamma_D^2}.
\]
材料因子的全半高宽为 $\Gamma_D$；只有 $\Gamma_D|\partial_\omega\log K_0(2\omega b/v)|\ll1$ 时，完整谱峰才近似继承这一线宽。

若 $\hbar\omega_p=9\,\mathrm{eV}$，静电极点为 $6.364\,\mathrm{eV}$。完整迟滞计算在 $v=0.70c$ 时的同步尺度却是
\[
\hbar\omega_{\rm sync}
=9\sqrt{\frac{1-0.70^2}{2-0.70^2}}\,\mathrm{eV}
\simeq5.230\,\mathrm{eV}.
\]
这两个数属于不同近似；下一个完整 Fresnel 例题中的约 $5.249\,\mathrm{eV}$ 谱峰不能由 $1+\varepsilon=0$ 直接解释。

当 $X=2\omega b/v\gg1$ 时，用 $\cosh u=1+u^2/2+O(u^4)$ 作 Laplace 展开，得到
\[
K_0(X)=\sqrt{\frac\pi{2X}}\ee^{-X}[1+O(X^{-1})].
\]
故增加距离 $\Delta b$ 主要压低高频信号 $\ee^{-2\omega\Delta b/v}$。束流有位置分布时，应平均这条概率谱，而不是把平均距离代入指数。
\end{workedexample}
